\appendix
\section{Short derivations}

\subsection{The window law}

For a fixed suffix $s$, write $y_-(s)$ and $y_+(s)$ for the infimum and
supremum of its compatible outputs.  Any suffix predictor returns one number
on this set, so values arbitrarily close to the two endpoints force an error
of at least
$[y_+(s)-y_-(s)]/2$.  The midpoint
$[y_+(s)+y_-(s)]/2$ reaches this bound.  Taking the supremum over suffixes proves
Eq.~\eqref{eq:window-law}.  The same lower bound holds for a randomized output
$A$, because
$|y_+-y_-|\leq\mathbb E|y_+-A|+\mathbb E|y_--A|$.  For a balanced pair,
\begin{equation}
 \frac12\sum_{b=0}^1\mathbb E(y_b-A)^2
 =\frac{(y_1-y_0)^2}{4}
  +\mathbb E\!\left[A-\frac{y_0+y_1}{2}\right]^2 .
\end{equation}

\subsection{The protected singlet}

If every allowed suffix channel satisfies $\Phi_a^\dagger(P_-)=P_-$, then any
composition of those channels does too.  The singlet population difference is
therefore unchanged by every finite allowed suffix.  The old pulse loads it
because
\begin{equation}
 \begin{aligned}
 R_x^{(1)}(\theta)|00\rangle
 &=\cos(\theta/2)|00\rangle-i\sin(\theta/2)|10\rangle,\\
 |\langle\psi^-|10\rangle|^2&=\frac12 .
 \end{aligned}
\end{equation}
This gives $\Delta p_-=\tfrac12\sin^2(\theta/2)$, while
$P_-=(II-XX-YY-ZZ)/4$ makes the population measurable from the stated Pauli
outputs.

For imperfect protection, let
$\varepsilon_{\rm prot}=\sup_a\|\Phi_a^\dagger(P_-)-P_-\|_\infty$.  Each
allowed step can reduce the branch gap by at most
$2\varepsilon_{\rm prot}$, because the trace norm of the difference between
two states is at most two.  Summing over $d$ steps gives the readable
finite-horizon bound
\begin{equation}
 |\Delta p_d|\geq
 \max\{0,|\Delta p_0|-2d\varepsilon_{\rm prot}\}.
\end{equation}

\subsection{The reset-QELM bound}

Let $\widetilde\rho_t^{(d)}$ be the state obtained by replaying the final $d$
channels from one fixed reset state.  Repeated contraction gives
\begin{equation}
 D_{\rm tr}(\rho_t,\widetilde\rho_t^{(d)})\leq\eta^d,
 \qquad
 |Y_t-\operatorname{Tr}(O\widetilde\rho_t^{(d)})|\leq2M\eta^d .
\end{equation}

For each window coordinate $x_i$, prepare $r$ qubits in
$\sqrt{1-x_i}|0\rangle+\sqrt{x_i}|1\rangle$ and record the Hamming weight
$K_i$.  Then $K_i\sim\operatorname{Binomial}(r,x_i)$.  If $g_d$ is the replay
output, a linear readout with coefficients $g_d(\boldsymbol k/r)$ returns the
Bernstein polynomial $\mathbb E[g_d(\boldsymbol K/r)]$.  Input Lipschitzness
and contraction give
\begin{equation}
 |g_d(\boldsymbol x)-g_d(\boldsymbol y)|
 \leq2ML_u\sum_{i=1}^d\eta^{d-i}|x_i-y_i|.
\end{equation}
Since
$\mathbb E|K_i/r-x_i|\leq1/(2\sqrt r)$, the Bernstein error is at most
$ML_u(1-\eta^d)/[(1-\eta)\sqrt r]$.  Finally, each circuit shot produces a
score in $[-M,M]$.  Hoeffding's inequality bounds the error of the mean of $Q$
scores by
$M\sqrt{2\log(2/\alpha)/Q}$ with probability at least $1-\alpha$.  Adding the
three errors proves Eq.~\eqref{eq:qrc-qelm-bound}.

For a general continuous fading-memory filter, choose a finite window whose
truncation error is below half the target tolerance.  The remaining map is a
continuous function on a compact finite-dimensional domain, so Bernstein
approximation supplies a finite reset QELM for the other half.  This proves the
stated general corollary without assuming exponential contraction.
